% Test Case 01: CRRA Relative Risk Aversion (Sanity Check)
% Difficulty: Easy — algebraic simplification only
%
% Claim: Under CRRA utility u(c) = c^{1-γ}/(1-γ), the coefficient of
% relative risk aversion is constant and equal to γ.
%
% The coefficient of relative risk aversion is defined as:
%   RRA(c) = -c · u''(c) / u'(c)
%
% Given u(c) = c^{1-γ}/(1-γ):
%   u'(c)  = c^{-γ}
%   u''(c) = -γ · c^{-γ-1}
%
% Therefore:
%   RRA(c) = -c · (-γ · c^{-γ-1}) / c^{-γ}
%           = γ · c^{-γ} / c^{-γ}
%           = γ

\begin{claim}
Under CRRA utility $u(c) = \frac{c^{1-\gamma}}{1-\gamma}$ where $\gamma > 0$ and
$\gamma \neq 1$, the coefficient of relative risk aversion is constant and equal to
$\gamma$. Formally, for all $c > 0$:
\[
  -c \cdot \frac{u''(c)}{u'(c)} = \gamma
\]
\end{claim}
